\documentclass{amsart}
% For dealing with references we use the comment environment
\usepackage{verbatim}
\newenvironment{reference}{\comment}{\endcomment}
%\newenvironment{reference}{}{}
\newenvironment{slogan}{\comment}{\endcomment}
\newenvironment{history}{\comment}{\endcomment}

% For commutative diagrams we use Xy-pic
\usepackage[all]{xy}

% We use 2cell for 2-commutative diagrams.
\xyoption{2cell}
\UseAllTwocells

% We use multicol for the list of chapters between chapters
\usepackage{multicol}

% This is generally recommended for better output
\usepackage{lmodern}
\usepackage[T1]{fontenc}

% For cross-file-references

% Package for hypertext links:
\usepackage{hyperref}

% For any local file, say "hello.tex" you want to link to please

% Theorem environments.
%
\theoremstyle{plain}
\newtheorem{theorem}[subsection]{Theorem}
\newtheorem{proposition}[subsection]{Proposition}
\newtheorem{lemma}[subsection]{Lemma}

\theoremstyle{definition}
\newtheorem{definition}[subsection]{Definition}
\newtheorem{example}[subsection]{Example}
\newtheorem{exercise}[subsection]{Exercise}
\newtheorem{situation}[subsection]{Situation}

\theoremstyle{remark}
\newtheorem{remark}[subsection]{Remark}
\newtheorem{remarks}[subsection]{Remarks}

\numberwithin{equation}{subsection}

% Macros
%
\def\lim{\mathop{\mathrm{lim}}\nolimits}
\def\colim{\mathop{\mathrm{colim}}\nolimits}
\def\Spec{\mathop{\mathrm{Spec}}}
\def\Hom{\mathop{\mathrm{Hom}}\nolimits}
\def\Ext{\mathop{\mathrm{Ext}}\nolimits}
\def\SheafHom{\mathop{\mathcal{H}\!\mathit{om}}\nolimits}
\def\SheafExt{\mathop{\mathcal{E}\!\mathit{xt}}\nolimits}
\def\Sch{\mathit{Sch}}
\def\Mor{\mathop{\mathrm{Mor}}\nolimits}
\def\Ob{\mathop{\mathrm{Ob}}\nolimits}
\def\Sh{\mathop{\mathit{Sh}}\nolimits}
\def\NL{\mathop{N\!L}\nolimits}
\def\CH{\mathop{\mathrm{CH}}\nolimits}
\def\proetale{{pro\text{-}\acute{e}tale}}
\def\etale{{\acute{e}tale}}
\def\QCoh{\mathit{QCoh}}
\def\Ker{\mathop{\mathrm{Ker}}}
\def\Im{\mathop{\mathrm{Im}}}
\def\Coker{\mathop{\mathrm{Coker}}}
\def\Coim{\mathop{\mathrm{Coim}}}

% Boxtimes
%
\DeclareMathSymbol{\boxtimes}{\mathbin}{AMSa}{"02}

%
% Macros for moduli stacks/spaces
%
\def\QCohstack{\mathcal{QC}\!\mathit{oh}}
\def\Cohstack{\mathcal{C}\!\mathit{oh}}
\def\Spacesstack{\mathcal{S}\!\mathit{paces}}
\def\Quotfunctor{\mathrm{Quot}}
\def\Hilbfunctor{\mathrm{Hilb}}
\def\Curvesstack{\mathcal{C}\!\mathit{urves}}
\def\Polarizedstack{\mathcal{P}\!\mathit{olarized}}
\def\Complexesstack{\mathcal{C}\!\mathit{omplexes}}
% \Pic is the operator that assigns to X its picard group, usage \Pic(X)
% \Picardstack_{X/B} denotes the Picard stack of X over B
% \Picardfunctor_{X/B} denotes the Picard functor of X over B
\def\Pic{\mathop{\mathrm{Pic}}\nolimits}
\def\Picardstack{\mathcal{P}\!\mathit{ic}}
\def\Picardfunctor{\mathrm{Pic}}
\def\Deformationcategory{\mathcal{D}\!\mathit{ef}}

\setlength{\emergencystretch}{3em}
\begin{document}
\title[Stacks Project, Tag 02O5]{372 --- Proper pushforward preserves coherent cohomology over locally Noetherian bases}
\maketitle
\noindent Copyright (C) 2005--2025 Johan de Jong. Stacks Project authors. Source: \href{https://stacks.math.columbia.edu/tag/02O5}{Tag 02O5}.

\noindent Editorial difficulty note (not author text). Difficulty H2: The proof must reduce arbitrary coherent modules to suitable generically rank-one sheaves whose direct images are controlled. A Chow modification and a twist ample for two maps create such test objects with vanishing higher direct images, and relative Leray plus coherent devissage complete the finiteness argument..

\noindent Editorial provenance note (not author text). History: excerpt from revision \texttt{a0444\allowbreak 6e57e\allowbreak c1fbc\allowbreak 252a8\allowbreak 71afc\allowbreak ec775\allowbreak 2fb28\allowbreak 07b14}, coherent.tex, lines 5107--5224. Reference presentation was adapted; original-author context and core support blocks were retained or added. The mathematical wording is unchanged. Licensed under GNU FDL 1.2 or later; no invariant sections or cover texts. See ../licenses/COPYING. Context blocks reproduce author definitions and supporting results; they are not additional counted problems.

\begin{proposition}
\label{proposition-proper-pushforward-coherent}
\begin{reference}
\href{https://stacks.math.columbia.edu/bibliography}{Bibliography: EGA (III Theorem 3.2.1)}
\end{reference}
Let $S$ be a locally Noetherian scheme.
Let $f : X \to S$ be a proper morphism.
Let $\mathcal{F}$ be a coherent $\mathcal{O}_X$-module.
Then $R^if_*\mathcal{F}$ is a coherent $\mathcal{O}_S$-module
for all $i \geq 0$.
\end{proposition}

\begin{proof}
Since the problem is local on $S$ we may assume that $S$ is
a Noetherian scheme. Since a proper morphism is of finite type
we see that in this case $X$ is a Noetherian scheme also.
Consider the property $\mathcal{P}$ of coherent sheaves
on $X$ defined by the rule
$$
\mathcal{P}(\mathcal{F}) \Leftrightarrow
R^pf_*\mathcal{F}\text{ is coherent for all }p \geq 0
$$
We are going to use the result of
Lemma \href{https://stacks.math.columbia.edu/tag/01YI}{lemma property (Tag 01YI)} to prove that
$\mathcal{P}$ holds for every coherent sheaf on $X$.

\medskip\noindent
Let
$$
0 \to \mathcal{F}_1 \to \mathcal{F}_2 \to \mathcal{F}_3 \to 0
$$
be a short exact sequence of coherent sheaves on $X$.
Consider the long exact sequence of higher direct images
$$
R^{p - 1}f_*\mathcal{F}_3 \to
R^pf_*\mathcal{F}_1 \to
R^pf_*\mathcal{F}_2 \to
R^pf_*\mathcal{F}_3 \to
R^{p + 1}f_*\mathcal{F}_1
$$
Then it is clear that if 2-out-of-3 of the sheaves $\mathcal{F}_i$
have property $\mathcal{P}$, then the higher direct images of the
third are sandwiched in this exact complex between two coherent
sheaves. Hence these higher direct images are also coherent by
Lemma \href{https://stacks.math.columbia.edu/tag/01Y0}{lemma coherent abelian Noetherian (Tag 01Y0)} and
\href{https://stacks.math.columbia.edu/tag/01Y1}{lemma coherent Noetherian quasi coherent sub quotient (Tag 01Y1)}.
Hence property $\mathcal{P}$ holds for the third as well.

\medskip\noindent
Let $Z \subset X$ be an integral closed subscheme.
We have to find a coherent sheaf $\mathcal{F}$ on $X$ whose support is
contained in $Z$, whose stalk at the generic point $\xi$ of $Z$ is a
$1$-dimensional vector space over $\kappa(\xi)$ such that $\mathcal{P}$
holds for $\mathcal{F}$. Denote $g = f|_Z : Z \to S$ the restriction of $f$.
Suppose we can find a coherent sheaf $\mathcal{G}$ on $Z$ such
that
(a) $\mathcal{G}_\xi$ is a $1$-dimensional vector space over $\kappa(\xi)$,
(b) $R^pg_*\mathcal{G} = 0$ for $p > 0$, and
(c) $g_*\mathcal{G}$ is coherent. Then we can consider
$\mathcal{F} = (Z \to X)_*\mathcal{G}$. As $Z \to X$ is a closed immersion
we see that $(Z \to X)_*\mathcal{G}$ is coherent on $X$
and $R^p(Z \to X)_*\mathcal{G} = 0$ for $p > 0$
(Lemma \href{https://stacks.math.columbia.edu/tag/01Y6}{lemma finite pushforward coherent (Tag 01Y6)}).
Hence by the relative Leray spectral sequence
(Cohomology, Lemma \href{https://stacks.math.columbia.edu/tag/01F6}{lemma relative Leray (Tag 01F6)})
we will have $R^pf_*\mathcal{F} = R^pg_*\mathcal{G} = 0$ for $p > 0$
and $f_*\mathcal{F} = g_*\mathcal{G}$ is coherent.
Finally $\mathcal{F}_\xi = ((Z \to X)_*\mathcal{G})_\xi = \mathcal{G}_\xi$
which verifies the condition on the stalk at $\xi$.
Hence everything depends on finding a coherent sheaf $\mathcal{G}$
on $Z$ which has properties (a), (b), and (c).

\medskip\noindent
We can apply Chow's Lemma \href{https://stacks.math.columbia.edu/tag/0200}{lemma chow Noetherian (Tag 0200)}
to the morphism $Z \to S$. Thus we get a diagram
$$
\xymatrix{
Z \ar[rd]_g & Z' \ar[d]^-{g'} \ar[l]^\pi \ar[r]_i & \mathbf{P}^m_S \ar[dl] \\
& S &
}
$$
as in the statement of Chow's lemma. Also, let $U \subset Z$ be
the dense open subscheme such that $\pi^{-1}(U) \to U$ is an isomorphism.
By the discussion in Remark \href{https://stacks.math.columbia.edu/tag/0201}{remark chow Noetherian (Tag 0201)} we see that
$i' = (i, \pi) : Z' \to \mathbf{P}^m_Z$ is
a closed immersion. Hence
$$
\mathcal{L} = i^*\mathcal{O}_{\mathbf{P}^m_S}(1) \cong
(i')^*\mathcal{O}_{\mathbf{P}^m_Z}(1)
$$
is $g'$-relatively ample and $\pi$-relatively ample (for example by
Morphisms, Lemma \href{https://stacks.math.columbia.edu/tag/02NR}{lemma characterize ample on finite type (Tag 02NR)}).
Hence by Lemma \href{https://stacks.math.columbia.edu/tag/02O1}{lemma kill by twisting (Tag 02O1)}
there exists an $n \geq 0$ such that
both $R^p\pi_*\mathcal{L}^{\otimes n} = 0$ for all $p > 0$ and
$R^p(g')_*\mathcal{L}^{\otimes n} = 0$ for all $p > 0$.
Set $\mathcal{G} = \pi_*\mathcal{L}^{\otimes n}$.
Property (a) holds because $\pi_*\mathcal{L}^{\otimes n}|_U$ is
an invertible sheaf (as $\pi^{-1}(U) \to U$ is an isomorphism).
Properties (b) and (c) hold because by the relative Leray
spectral sequence
(Cohomology, Lemma \href{https://stacks.math.columbia.edu/tag/01F6}{lemma relative Leray (Tag 01F6)})
we have
$$
E_2^{p, q} = R^pg_* R^q\pi_*\mathcal{L}^{\otimes n}
\Rightarrow
R^{p + q}(g')_*\mathcal{L}^{\otimes n}
$$
and by choice of $n$ the only nonzero terms in $E_2^{p, q}$ are
those with $q = 0$ and the only nonzero terms of
$R^{p + q}(g')_*\mathcal{L}^{\otimes n}$ are those with $p = q = 0$.
This implies that $R^pg_*\mathcal{G} = 0$ for $p > 0$ and
that $g_*\mathcal{G} = (g')_*\mathcal{L}^{\otimes n}$.
Finally, applying the previous
Lemma \href{https://stacks.math.columbia.edu/tag/02O4}{lemma locally projective pushforward (Tag 02O4)}
we see that $g_*\mathcal{G} = (g')_*\mathcal{L}^{\otimes n}$ is
coherent as desired.
\end{proof}
\end{document}
